\begin{figure}[ht]
\centering
\begin{tikzpicture}[x=1cm, y=1cm, font=\small,
  family/.style={draw=engblue, thick, rounded corners=4pt, fill=engblue!8,
                 minimum width=3.6cm, minimum height=1.0cm, align=center,
                 font=\bfseries\footnotesize},
  example/.style={font=\scriptsize, anchor=west, align=left}
]

\node[family] (thermal)   at (0,  4.5) {Thermal};
\node[family] (electrical) at (4.2, 4.5) {Electrical};
\node[family] (optical)   at (8.4, 4.5) {Optical};
\node[family] (mechanical) at (0,  1.5) {Mechanical};
\node[family] (chemical)  at (4.2, 1.5) {Chemical};
\node[family] (magnetic)  at (8.4, 1.5) {Magnetic};

% example lists
\node[example, anchor=north west] at (-1.8, 3.9)
  {thermocouple\\RTD (Pt100)\\thermistor\\IR pyrometer};
\node[example, anchor=north west] at (2.4, 3.9)
  {shunt, CT, PT\\Hall current\\charge amp.\\capacitive probe};
\node[example, anchor=north west] at (6.6, 3.9)
  {photodiode\\PMT\\CCD / CMOS\\fibre, interferometer};

\node[example, anchor=north west] at (-1.8, 0.9)
  {strain gauge\\piezo (PZT, PVDF)\\LVDT\\MEMS, pressure};
\node[example, anchor=north west] at (2.4, 0.9)
  {ISE (pH, Na$^+$)\\electrochemical\\optical / SPR\\metal-oxide gas};
\node[example, anchor=north west] at (6.6, 0.9)
  {Hall effect\\AMR / GMR / TMR\\fluxgate\\SQUID};

% center hub
\node[draw, thick, dashed, circle, fill=engred!10, minimum size=1.8cm,
      align=center, font=\footnotesize] (hub) at (4.2, 3.0)
      {sensor\\datasheet\\contract};

\foreach \src in {thermal, electrical, optical, mechanical, chemical, magnetic} {
  \draw[thin, enggray] (\src) -- (hub);
}

\end{tikzpicture}
\caption[The six transducer families and their working examples]{The
catalogue a working engineer should recognise before reading a
datasheet. Each family converts a physical input into an electrical
(or otherwise readable) output by a distinct physical principle.
The boundaries are working, not exclusive: an LVDT is mechanical
in the input quantity it senses but inductive in its readout
principle, and a MEMS accelerometer is mechanical in input and
capacitive or piezoresistive in readout. The datasheet contract at
the centre is what binds the device to a measurement, regardless of
family.}
\label{fig:vol01:ch05:transducer-families}
\end{figure}
